\documentclass[12pt]{article}
\usepackage[margin=1in]{geometry}
\usepackage[all]{xy}


\usepackage{amsmath,amsthm,amssymb,color,latexsym}
\usepackage{geometry}        
\geometry{letterpaper}  
\usepackage{boxedminipage}
\usepackage{graphicx}
\usepackage{hyperref}
\hypersetup{
    colorlinks=true,
    citecolor=blue,
    linkcolor=blue,
    filecolor=cyan,
    urlcolor=magenta,
}
\usepackage[capitalize,nameinlink]{cleveref}

\newtheorem{problem}{Problem}
\newtheorem{lemma}{Lemma}


\newenvironment{solution}[1][\it{Solution}]{\textbf{#1. } }{$\square$}
\newcommand{\pointsbox}[1]{%
  \colorbox{red!70}{\strut\textcolor{white}{\sffamily\bfseries\ \ #1\ \ }}%
}
\newcommand{\pagetag}[1]{%
  \colorbox{red!70}{\strut\textcolor{white}{\sffamily\bfseries\ \ #1\ \ }}%
}
\newcommand{\striptitle}[2]{%
  \noindent
  \begin{minipage}{0.07\linewidth}\centering \pagetag{#1}\end{minipage}%
  \hfill
  \begin{minipage}{0.89\linewidth}
    {\large\scshape #2}
  \end{minipage}\par\vspace{0.75em}
}
% --- Environments ---
\theoremstyle{plain}
\newtheorem{theorem}{Theorem}
\newtheorem{corollary}[theorem]{Corollary}
\theoremstyle{definition}
\newtheorem{definition}[theorem]{Definition}
\newtheorem{remark}[theorem]{Remark}

\newenvironment{boxedalgo}
  {\begin{center}\begin{boxedminipage}{1\linewidth}}
  {\end{boxedminipage}\end{center}}

\newenvironment{inlinealgo}[1]
  {\noindent\hrulefill\\\noindent{\bf {#1}}}
  {\hrulefill}


% --- Shortcuts ---
\newcommand{\R}{\mathbb{R}}
\newcommand{\inner}[2]{\left\langle #1,#2 \right\rangle}
\newcommand{\proj}{\mathrm{proj}}
\newcommand{\Id}{\mathbf{I}}
\newcommand{\ot}{\otimes}
\newcommand{\mc}{\mathcal}

\begin{document}
\noindent CS 578 Statistical Machine Learning (Fall 2026)\hfill Problem Set \# 3\\
\{\{Your name\}\} \hfill \textbf{Due: September 25th}

\noindent\hrulefill

\subsection*{Problem 1. SVM.}
In this problem, you will prove that the hard-SVM rule, namely,
$$
\operatornamewithlimits{argmax}_{(w,b):\|w\|=1} \min_{i\in [m]} |\langle w, x_i\rangle + b|\quad \text{s.t.}\quad \forall i,~y_i(\langle w, x_i\rangle + b)>0,
$$
is equivalent to the following formulation:
$$
\operatornamewithlimits{argmax}_{(w,b):\|w\|=1} \min_{i\in [m]} y_i(\langle w, x_i\rangle + b).
$$
\begin{enumerate}
    \item Define $\mathcal{G}:=\{(w,b):\forall i,~y_i(\langle w,x_i\rangle + b)>0\}.$ Show that 
    $$
        \operatornamewithlimits{argmax}_{(w,b):\|w\|=1} \min_{i\in [m]} y_i(\langle w, x_i\rangle + b)\in \mathcal{G}.
    $$
    \item Show that for any $(w,b)\in {\cal G}$,
    $$
        \min_{i\in [m]} y_i(\langle w,x_i\rangle + b)=\min_{i\in [m]} |\langle w,x_i\rangle + b|,
    $$
    and finish the proof of the equivalence.
\end{enumerate}
\noindent{\bfseries Solution.}\newline 
\   %%% <-- ERASE THIS LINE AND WRITE YOUR SOLUTION HERE
\newpage

\subsection*{Problem 2. Rademacher Complexity.}
In this problem, you will prove the Massart lemma:
\begin{lemma}[Massart lemma]
Let  $A = \{a_1,\dots, a_N\}$ be a finite set of vectors in $\mathbb{R}^m$. Define $\overline{a}=\frac{1}{N}\sum_i a_i$. Then, the Rademacher complexity of $A$ can be bounded by:
\begin{align*}
    {\rm R}(A)\leq \max_{i\in [N]} \|a_i-\overline{a}\|\cdot \frac{\sqrt{2\log N}}{m}.
\end{align*}
\end{lemma}

\end{document}
